% Starter Beamer deck for HEP talks. 16:9, Madrid theme (ships with TeX Live, no
% extra install). Copy next to a figures/ directory and edit.
% Build: latexmk -pdf -interaction=nonstopmode -halt-on-error talk.tex
\documentclass[aspectratio=169,12pt]{beamer}

\usetheme{Madrid}
\usecolortheme{default}
\setbeamertemplate{navigation symbols}{}
\setbeamertemplate{footline}[frame number]

\usepackage{graphicx}
\usepackage{booktabs}
\usepackage{amsmath}
\usepackage{siunitx}
\graphicspath{{figures/}}

% Frame titles state the takeaway; keep them to two lines at most.
\setbeamerfont{frametitle}{size=\large}

% One-line takeaway under a figure.
\newcommand{\takeaway}[1]{\par\vspace{10pt}\centering\small\textbf{#1}\par}

\title[Short title]{Full talk title}
\subtitle{Analysis / channel}
\author[Presenter]{Presenter Name}
\institute[OSU]{The Ohio State University}
\date{\today}

\begin{document}

\begin{frame}
  \titlepage
\end{frame}

% ---------------------------------------------------------------- Motivation
\section{Motivation}
\begin{frame}{State the physics goal as a takeaway}
  \begin{itemize}
    \item At most ~4 bullets, ~2 lines each
    \item Detail belongs in backup
  \end{itemize}
\end{frame}

% ------------------------------------------------------ Single-figure frame
\section{Results}
\begin{frame}{Title states the result, not the topic}
  \centering
  % Size by HEIGHT so plots with/without colorbars match (see slide-figures.md).
  \includegraphics[height=0.72\textheight]{example_plot.pdf}
  \takeaway{One-line message the plot supports}
\end{frame}

% ----------------------------------------------------- Two-figure frame
\begin{frame}{Compare two views side by side}
  \centering
  \includegraphics[height=0.62\textheight]{plot_a.pdf}\hfill
  \includegraphics[height=0.62\textheight]{plot_b.pdf}
  \takeaway{What the comparison shows}
\end{frame}

% ------------------------------------------------------------ Table frame
\begin{frame}{Cutflow (values from the analysis output, not typed by hand)}
  \centering
  \normalsize   % tables stay at body size; only go \small if a table truly won't fit
  % Free-text columns: use p{<width>} so they wrap (an l column overruns the slide).
  \begin{tabular}{lrr}
    \toprule
    Selection & Data & MC \\
    \midrule
    % \input{tables/cutflow_body.tex}   % generate from the analysis output
    Placeholder & -- & -- \\
    \bottomrule
  \end{tabular}
\end{frame}

% ---------------------------------------------------------------- Summary
\begin{frame}{Summary and next steps}
  \begin{itemize}
    \item Key result, with its uncertainty
    \item Concrete next step
  \end{itemize}
\end{frame}

% ---------------------------------------------------------------- Backup
\appendix
\section{Backup}
\begin{frame}[noframenumbering]{Backup: extra validation}
  \centering
  \includegraphics[height=0.72\textheight]{backup_plot.pdf}
\end{frame}

\end{document}
